\bookchapter{الفصل 25: معالجة الهواتف المستعملة والتراث الرقمي}

\chapterintro{مرحبًا يا رفاق، بعد “التدريب” في أكثر من عشرين فصلًا سابقة، أعتقد أنك تطورت من “مبتدئ رقمي” إلى مستخدم هاتف واثق تمامًا. لكن الحياة التقنية لها دورة — سيتحول الحبيب الجديد قريبًا إلى حبيب قديم. عندما تستعد لإحالة “الرفيق القديم” إلى التقاعد، هل تعرف كيف تودعه بشكل لائق وآمن؟}

\booksection{أولًا: “الموت الرحيم” للهاتف القديم — دليل البيع أو الإهداء الآمن}

المبدأ الأساسي: تنظيف البيانات ليس “حذفاً” بل “استبدالاً”. كثيرون يظنون أن سحب الصور إلى سلة المحذوفات وإفراغ الألبوم يكفي، لكن بعد بيع الهاتف، يستخدم المالك الجديد برنامج استعادة بيانات فيسترجع الصور القديمة وسجلات الدردشة وصور الهوية — هذا “مشهد موت اجتماعي” يتكرر كل عام.

\begin{enumerate}
  \item \textbf{: اضبط مدة، وعند حلولها يُرسل تلقائيًا إلى جهة اتصال محددة أو يُحذف الحساب.}: أكمل آخر نسخة احتياطية كاملة، وسجل الخروج من جميع الحسابات (iCloud وHuawei/Xiaomi/Samsung وGoogle)، وأغلق وظيفة "Find My Phone" — وإلا سيحجب قفل التنشيط المستخدم التالي.
  \item \textbf{: الطريقة الأكثر موثوقية، نزّل البيانات الثمينة بانتظام إلى قرص صلب محلي.}: "الاستعادة للمصنع" في الهواتف الحديثة تتضمن عادةً مسحاً مشفّراً — تدمر مفتاح التشفير نفسه، وبمجرد فقدان المفتاح تتحول بيانات القرص بالكامل إلى بيانات غير مفهومة.
  \item \textbf{"التمويه" الأخير (الأجهزة القديمة فقط)}: الهواتف الحديثة (iOS 8+ وAndroid 6+) مشفّرة بالكامل افتراضياً (FBE)، والاستعادة للمصنع تدمر مفتاح التشفير، والبيانات لا يمكن استعادتها نظرياً،\textbf{لا حاجة لهذه الخطوة}. فقط للأجهزة القديمة جداً غير المشفّرة (قبل 2015)، قد يحتاج ملء التخزين بملفات غير مهمة بعد الاستعادة للمصنع، ثم تنفيذ استعادة للمصنع مرة أخرى.
\end{enumerate}

\textbf{جدول مرجعي سريع لمسار مسح البيانات لكل علامة}: رغم أن المبدأ واحد، إلا أن كل علامة تخفي “الاستعادة للمصنع” في مكان مختلف قليلاً، والاختلافات الرئيسية كالتالي.

\begin{center}
\begin{tabularx}{\linewidth}{|X|X|X|}
\hline
\textbf{العلامة} & \textbf{مدخل المسح} & \textbf{الاختلافات والملاحظات الرئيسية} \tabularnewline
\hline
Huawei & الإعدادات > النظام والتحديث > إعادة تعيين > الاستعادة للمصنع & يجب الخروج من حساب Huawei وإغلاق "Find My Phone" أولاً، وإلا يُفعّل قفل التنشيط؛ قد يتطلب التحقق بكلمة المرور في الوضع النقي. \tabularnewline
Xiaomi / Redmi & الإعدادات > جهازي > الاستعادة للمصنع & في الأسفل خيارا "تهيئة SD card الافتراضية" و"مسح جميع المحتويات والإعدادات"، يجب تحديد كليهما؛ أغلق "Find Device" قبل الخروج من حساب Xiaomi. \tabularnewline
OPPO / OnePlus & الإعدادات > إعدادات أخرى > استعادة الهاتف > مسح جميع البيانات بالكامل & ColorOS يوفر وضع "المسح الكامل"، يغطي التخزين عدة مرات؛ يتطلب إدخال كلمة مرور قفل الشاشة للتأكيد. \tabularnewline
vivo & الإعدادات > إدارة النظام > النسخ الاحتياطي وإعادة التعيين > مسح جميع البيانات & تحديد "تهيئة تخزين الهاتف" يمسح منطقة SD card؛ يجب الخروج من حساب vivo وإغلاق "Find Phone". \tabularnewline
Samsung & الإعدادات > الإدارة العامة > إعادة تعيين > الاستعادة للمصنع & يجب إزالة حساب Samsung وإغلاق "Find My Phone" أولاً؛ بعض الموديلات لها خيار "تغطية تلقائية"، يُنصح بتفعيله. \tabularnewline
iPhone & الإعدادات > عام > نقل أو استعادة iPhone > مسح جميع المحتويات والإعدادات & يجب الخروج من Apple ID أولاً، وإلا يُحجب المستخدم التالي بقفل تنشيط iCloud؛ iOS 15+ يدعم "الاحتفاظ ببيانات iCloud عند الخروج" للمراجعة لاحقاً. \tabularnewline
\hline
\end{tabularx}
\end{center}

\begin{tcbwarn}{ثلاثة أشياء يسهل تجاهلها أكثر}
\begin{itemize}
  \item \textbf{عدم الخروج من الحساب}: الهاتف يُقفل عن بُعد بعد بيعه، ولا تستطيع الخروج أيضاً، وفي النهاية تضطر لزيارة خدمة ما بعد البيع الرسمية بالهوية لفتح القفل.
  \item \textbf{عدم إخراج SD card}: الاستعادة للمصنع لا تهيئ الصور على SD card الخارجية، تذكر إخراجها يدوياً أو تهيئتها بشكل منفصل.
  \item \textbf{عدم فصل فلاش USB/أجهزة OTG}: إذا كانت الأجهزة الخارجية متصلة أثناء المسح، قد تُمسح بياناتها بالخطأ، أو قد يتوقف المسح.
\end{itemize}

\end{tcbwarn}

\booksection{ثانيًا: التراث الرقمي — أين تضع روحك الاجتماعية؟}

البيانات في الهاتف لم تعد “افتراضية” منذ زمن. تحمل أموالك وعلاقاتك الاجتماعية وإبداعاتك، بل و“هويتك الرقمية”.\textbf{بعد وفاة الشخص، كيف يجب التعامل مع هذه الأصول؟}هذا لم يعد سؤالاً مجرداً، بل معضلة حقيقية تواجهها أسر كثيرة — الأقارب لا يعرفون كلمات المرور، والمنصات لا تسمح بالميراث، والعملات المشفرة تُفقد للأبد.

\begin{tcbinfo}{الأساس القانوني: المادة 49 من قانون حماية المعلومات الشخصية}
ينص قانون حماية المعلومات الشخصية لجمهورية الصين الشعبية الذي دخل حيز التنفيذ 2021 في المادة 49 بوضوح:\textbf{عند وفاة شخص طبيعي، يجوز لأقاربه من الدرجة الأولى لمصلحته المشروعة والعادلة ممارسة الحقوق المنصوص عليها في هذا الفصل بشأن معلومات المتوفى الشخصية، مثل الاطلاع والنسخ والتصحيح والحذف؛ إلا إذا كان للمتوفى ترتيب آخر قبل وفاته.}

هذا يعني:

\begin{itemize}
  \item \textbf{للأقارب أساس قانوني}للتقدم للمنصات بطلب معالجة حسابات المتوفى ومعلوماته الشخصية.
  \item "المصلحة المشروعة والعادلة" تشمل عادةً معالجة الميراث وإغلاق الحسابات وحذف الخصوصية الحساسة وغيرها.
  \item إذا كان للمتوفى وصية أو اتفاق بترتيب آخر قبل وفاته (مثل تعيين جهة اتصال الميراث)، يُعمل بذلك.
  \item لا تستطيع المنصات الرفض المطلق للتعاون بحجة "خصوصية المستخدم"، لكن الإجراءات تختلف من منصة لأخرى.
\end{itemize}

\end{tcbinfo}

\booksubsection{WeChat وAlipay: الحساب لا يُورّث، لكن الأموال نعم}

هذان أكثر تطبيقين شيوعاً في الصين، ولكل منهما قواعد مختلفة:

\begin{itemize}
  \item \textbf{حساب WeChat}: وفقاً لـ "اتفاقية ترخيص واستخدام برنامج Tencent WeChat"، ملكية الحساب تعود لـ Tencent، والمستخدم يتمتع بحق الاستخدام فقط،\textbf{الحساب نفسه لا يمكن توريثه أو تحويله}. لكن الأموال في الحساب\textbf{رصيد WeChat Pay وأموال Liqiantong}تعتبر ممتلكات شخصية، ويمكن للأقارب التقدم إلى TenPay بشهادة الوفاة وشهادة القرابة وشهادة الميراث لسحب أو تحويل الأموال.
  \item \textbf{Alipay}: الحساب أيضاً لا يمكن توريثه مباشرة، لكن رصيد Alipay وYu'e Bao وLicaitong وغيرها\textbf{تعتبر أموالاً قابلة للتوريث}. يمكن للورثة المعالجة عبر "خط خدمة العملاء" في Alipay أو تقديم شهادة الوفاة وشهادة الميراث شخصياً؛ المنتجات المالية في Ant Fortune تحتاج إجراءات شركة الصناديق.
  \item \textbf{سجلات الدردشة}: قانونياً تعتبر خصوصية شخصية، ولن تقدمها المنصة للأهل بشكل استباقي. إذا تعلق الأمر بأدلة الميراث، يمكن طلبها عبر الدعوى القضائية أو التوثيق.
\end{itemize}

\begin{tcbwarn}{أشياء يجب ألا تفعلها أبداً}
\begin{itemize}
  \item \textbf{لا تسلم كلمة مرور حساب المتوفى للأسرة لـ "الاستخدام المباشر"}— هذا يخالف اتفاقية المنصة، وعند تجميد الأموال أو تفعيل الرقابة المخاطر يكون الإثبات أصعب.
  \item \textbf{لا تحاول استخدام هاتف المتوفى لاستلام رمز التحقق و"سحب الأموال"}— قد يُعتبر استيلاءً، مما يخالف القانون.
  \item \textbf{لا تتأخر}: بعض الحسابات تُلغى تلقائياً بعد عدم تسجيل الدخول لفترة طويلة، واسترجاع الأموال سيكون أصعب.
\end{itemize}

\end{tcbwarn}

\booksubsection{“جهة اتصال الميراث” من Apple: التعيين المسبق يضاعف الكفاءة}

منذ iOS 15.2 أطلقت Apple وظيفة “Legacy Contact”، تتيح للمستخدم تعيين جهة اتصال موثوقة مسبقاً، وبعد الوفاة يمكن لهذه الجهة بـ\textbf{شهادة وفاة + مفتاح الوصول لجهة اتصال الميراث}التقدم للوصول إلى بيانات iCloud للمتوفى (الصور والملاحظات وجهات الاتصال وغيرها). خطوات العملية:

\begin{enumerate}
  \item على iPhone ادخل\textbf{[الإعدادات] > [اسمك] > [تسجيل الدخول والأمان] > [جهة اتصال الميراث]}.
  \item اضغط "إضافة جهة اتصال الميراث"، واستخدم Face ID لتأكيد الهوية.
  \item اختر جهة الاتصال (من دليل العناوين)، سيولّد النظام\textbf{مفتاح الوصول}.
  \item سلّم مفتاح الوصول للطرف الآخر عبر رسالة نصية أو بريد إلكتروني أو نسخة ورقية مطبوعة، وذكّره بالحفظ الجيد.
  \item يمكن إضافة عدة جهات اتصال كنسخ احتياطي، لتجنب فشل النقطة الواحدة.
\end{enumerate}

لا تحتاج الجهة المعيّنة لجهاز Apple، يمكنها التقدم على موقع Apple الرسمي بالمفتاح وشهادة الوفاة. تكتمل العملية عادةً خلال 7 أيام.

\booksubsection{“مدير الحساب غير النشط” من Google: معالجة آلية لما بعد الوفاة}

يقدم Google “مدير الحساب غير النشط (Inactive Account Manager)”، يمكنه تنفيذ عمليات مُعدة مسبقاً عند عدم تسجيل المستخدم الدخول لفترة طويلة — هذه أكثر آليات الميراث الرقمي اكتمالاً حالياً.

\begin{enumerate}
  \item املأ اسم المستودع (Repository name)، مثل\textbf{myaccount.google.com/inactive}سجّل الدخول لحساب Google.
  \item اضبط مهلة الانقطاع: 3 / 6 / 12 / 18 شهراً من عدم تسجيل الدخول للتشغيل.
  \item أضف جهات اتصال موثوقة (حتى 10 أشخاص)، واضبط البيانات التي يمكن لكل واحد الوصول إليها (Gmail والصور وDrive وغيرها).
  \item اختياري: ضبط رد تلقائي على البريد، يخبر المُراسلين بأنك لم تعد قادراً على المعالجة.
  \item الاختيار الأخير: بعد انتهاء المدة هل\textbf{تشارك البيانات مع جهات الاتصال}أم\textbf{تحذف الحساب مباشرة}.
  \item اضبط رقم هاتف كتحقق احتياطي، لتجنب التشغيل الخاطئ.
\end{enumerate}

\booksubsection{حسابات التواصل الاجتماعي المحلية: التحويل التذكاري والإلغاء}

تتعامل منصات التواصل الاجتماعي المحلية الرئيسية مع حسابات المتوفى بطريقتين:\textbf{“التحويل التذكاري” (الاحتفاظ بالحساب ليتذكره الأهل والأصدقاء)}و\textbf{“الإلغاء” (الحذف الكامل)}. يمكن للأهل التقدم للمنصة بشهادة الوفاة وشهادة القرابة وهوية مقدم الطلب.

\begin{itemize}
  \item \textbf{Weibo}: يدعم "التحويل التذكاري" للحساب، يقدمه الأهل بشهادة الوفاة، يُقفل الحساب ولا يمكن تسجيل الدخول إليه لكن يُحتفظ بالمحتوى، وتعرض الصفحة الرئيسية "هذا الحساب تحول إلى حساب تذكاري". يمكن أيضاً طلب الإلغاء، ويتطلب تقديم شهادة ميراث.
  \item \textbf{Douyin}: عبر "أنا - الإعدادات - التعليقات والمساعدة - الحساب والأمان - إلغاء الحساب" أو خدمة العملاء البشرية، قدم شهادة الوفاة وشهادة القرابة وهوية مقدم الطلب، وبعد الموافقة يمكن الإلغاء أو التحويل للحالة التذكارية.
  \item \textbf{Bilibili}: قدم عبر قناة خدمة العملاء شهادة المتوفى وشهادة القرابة، يمكن طلب إلغاء الحساب أو الاحتفاظ به كصفحة تذكارية؛ معالجة مقالات UP التاريخية تحتاج مراعاة إعلان حقوق النشر.
  \item \textbf{Xiaohongshu (RED)}: عبر تظلم خدمة العملاء الرسمي، قدم شهادة الوفاة وشهادة القرابة، يمكن طلب إلغاء الحساب وحذف المحتوى.
\end{itemize}

\begin{tcbinfo}{تذكير رئيسي}
\begin{itemize}
  \item قواعد كل منصة تتحدث، راجع أحدث إرشادات خدمة العملاء الرسمية قبل المعالجة.
  \item يُنصح بمسح جميع المستندات ضوئياً كنسخة احتياطية قبل التقديم، مع حفظ الأصل بعناية.
  \item عند المعالجة لكل المنصات في آن واحد، يمكن طلب\textbf{شهادة ميراث}واحدة موحدة، توفيراً لتكرار الإجراءات.
\end{itemize}

\end{tcbinfo}

\booksubsection{معضلة توريث الأصول الرقمية}

إلى جانب حسابات التواصل الاجتماعي، يمتلك المزيد من الناس\textbf{أصولاً رقمية ذات قيمة حقيقية}وتوريثها أكثر تعقيداً:

\begin{itemize}
  \item \textbf{العملات المشفرة (Bitcoin وEthereum وغيرها)}: بناءً على مبدأ "المفتاح الخاص هو الملكية"،\textbf{بدون المفتاح الخاص تختفي الأصول}. لا يستطيع الورثة "التظلم" على السلسلة عبر مستندات قانونية. يُنصح: تخزين كلمات الاستعادة/المفاتيح الخاصة مقسّمة في عدة مواقع آمنة مادية (مثل خزنة البنك + خزنة المنزل)، وذكر طريقة الوصول في الوصية.
  \item \textbf{حسابات منصات التداول}: مثل الحسابات البنكية، يمكن التقدم للمنصة (مثل Binance وOKX) بشهادة الميراث للمعالجة، لكن إجراءات المنصات عبر الحدود طويلة.
  \item \textbf{أسماء النطاقات}: تعتبر مالاً قابلاً للتوريث، يقدم الورثة شهادة الميراث لمسجل النطاق (مثل AliCloud وGoDaddy) لإتمام التحويل.
  \item \textbf{الملكية الفكرية (المقالات الأصلية والفيديوهات وبراءات الاختراع وحقوق نشر البرمجيات)}: الحقوق المالية في حق المؤلف قابلة للتوريث (الحقوق المعنوية لا تورّث)، وحق الإيرادات يتمتع به الورثة حتى 50 عاماً بعد وفاة المؤلف.
  \item \textbf{التخزين السحابي وأقراص الشبكة}: بيانات المحتوى تعود للفرد وقابلة للتوريث؛ لكن بعض اتفاقيات المنصات تنص على عدم إمكانية تحويل الحساب نفسه، يحتاج التنزيل أولاً ثم الإلغاء.
\end{itemize}

\textbf{السعر الحقيقي لـ “اختفاء” العملات المشفرة}\par

وفقاً لتقرير Chainalysis، هناك عالمياً\textbf{أكثر من 3 ملايين Bitcoin}لا يمكن استعادتها أبداً بسبب فقدان المفاتيح الخاصة، أي ما يعادل سدس الكمية الحالية. إذا كنت تمتلك أصولاً مشفّرة، احرص على إخبار شخص واحد على الأقل تثق به تماماً بـ\textbf{طريق الاستعادة}— وإلا ستبقى هذه الثروة نائمة على البلوكتشين إلى الأبد.

\booksubsection{نموذج “وصية رقمية”}

وصية الممتلكات التقليدية لا تغطي الأصول الرقمية. يُنصح بتحديث ملحق “وصية رقمية” سنوياً، وحفظه مع الوصية الرسمية. النموذج التالي للإرشاد:

\textbf{ملحق معالجة الأصول الرقمية الشخصية (نموذج)}\par

\textbf{القسم الأول: قائمة الحسابات}(سجل اسم الحساب والمنصة فقط،\textbf{لا تكتب كلمات المرور})

\begin{itemize}
  \item التواصل الاجتماعي: WeChat (الحساب xxx) وWeibo وDouyin وBilibili وXiaohongshu وغيرها
  \item المدفوعات: Alipay وWeChat Pay وتطبيقات البنوك
  \item السحابة: iCloud وHuawei Cloud وXiaomi Cloud وحساب Google وحساب Microsoft
  \item البريد: QQ Mail وNetEase Mail وGmail وغيرها
  \item أخرى: أقراص الشبكة والخدمات المشتركة وحسابات الألعاب وغيرها
\end{itemize}

\textbf{القسم الثاني: خطة حفظ كلمات المرور}

\begin{itemize}
  \item طريقة حفظ كلمة المرور الرئيسية: استخدام مدير كلمات المرور (مثل Bitwarden و1Password)، كتابة كلمة المرور الرئيسية على ورقة، حفظها في خزنة المنزل.
  \item مفاتيح الاستعادة/كلمات الاستعادة: تخزين مقسّم، موقعان ماديان على الأقل (مثل خزنة البنك + خزنة المنزل).
  \item الشخص المخول بالاستعلام: الاسم ورقم الهوية والعلاقة بالمستخدم ووسيلة الاتصال.
\end{itemize}

\textbf{القسم الثالث: نية المعالجة}

\begin{itemize}
  \item هل تريد "التحويل التذكاري" أم "الإلغاء" لكل نوع حساب.
  \item هل تم تعيين "جهة اتصال الميراث" أو "مدير الحساب غير النشط" لبيانات iCloud / Google.
  \item من يرث الأصول المشفرة، وشرح طريق الحصول على المفتاح الخاص.
  \item من يتمتع بإيرادات أعمال حقوق النشر.
\end{itemize}

\textbf{القسم الرابع: جهات الاتصال والمنصات}

\begin{itemize}
  \item اسم منفذ الثقة ووسيلة الاتصال وجهة الاتصال الاحتياطية.
  \item مداخل خدمة العملاء للمنصات وقائمة المستندات المطلوبة.
  \item تاريخ تحديث هذا الملحق: \_\_\_\_/\_\_/\_\_، تاريخ المراجعة التالية: \_\_\_\_/\_\_/\_\_.
\end{itemize}

\begin{tcbwarn}{تلميح قانوني حول "كلمات المرور"}
كتابة كلمات المرور مباشرة في الوصية ينطوي على خطر التسريب، وبعد إعلان الوصية يمكن لجميع الورثة رؤيتها.\textbf{الطريقة الأكثر أماناً: اكتب في الوصية “قائمة الحسابات + نية المعالجة + موقع حفظ كلمات المرور” فقط، ومرر كلمة المرور الحقيقية عبر كلمة المرور الرئيسية لمدير كلمات المرور بشكل منفصل لمنفذ الثقة}. هكذا يحقق المتطلبات القانونية ويخفض خطر تسريب كلمات المرور.

\end{tcbwarn}

\booksection{ثالثًا: الرحلة الأخيرة الصديقة للبيئة}

الهاتف الواحد يحوي عشرات المعادن كالذهب والفضة والنحاس والعناصر النادرة، ورميه عشوائياً يلوث التربة ويهدر الموارد. هناك ثلاثة قنوات إعادة تدوير صحيحة، لكل منها تركيز مختلف:

\begin{center}
\begin{tabularx}{\linewidth}{|X|X|X|X|}
\hline
\textbf{القناة} & \textbf{سعر الإعادة} & \textbf{أمان البيانات} & \textbf{سيناريوهات الاستخدام} \tabularnewline
\hline
الإعادة بجهاز قديم رسمي (Huawei/Apple/Xiaomi وغيرها) & متوسط، يُضاف إليه قسائم دعم غالباً & الأعلى، إجراءات رسمية موحدة & الأكثر جدوى عند التخطيط لتبديل هاتف نفس العلامة \tabularnewline
منصات إعادة التدوير الخارجية (Aihuishou وZhuanzhuan وHuishoubao وغيرها) & أعلى قليلاً، يمكن مقارنة الأسعار & عالٍ، يوجد تقرير مسح بيانات & لمن يريد تحويل نقدي مباشر دون الارتباط بعلامة \tabularnewline
التبرع العام (مثل مدارس الريف والمؤسسات الخيرية) & لا عائد اقتصادي & يجب المسح الكامل ذاتياً & جهاز لا يزال صالحاً للاستخدام، تريد فعل خير \tabularnewline
استغلال المتبقي (الإبقاء والتعديل) & توفير ثمن جهاز جديد & المالك الأصلي يستمر في الاستخدام، بلا مخاطر & أداء الجهاز لا يزال جيداً، هناك حاجة للتعديل في المنزل \tabularnewline
\hline
\end{tabularx}
\end{center}

\begin{itemize}
  \item \textbf{اختر محطته التالية: إعادة التدوير الرسمي، أو منصة معتمدة، أو إهداء.}: الطريقة الأكثر توصية، إجراءات موحدة، وأمان البيانات مضمون. مواقع Apple وHuawei وXiaomi وغيرها لها مدخل "trade-in"، وبعد التقييم يمكن خصمه من ثمن الجهاز الجديد.
  \item \textbf{وداعًا لطقوس "إعادة الضبط" عند التخلي عن الهاتف القديم، والتخطيط "للمصير الرقمي" لأصولك الرقمية — هذا ليس مجرد نهاية حسنة، بل أشبه برسالة مكتوبة للمستقبل. دعنا اليوم نطلق العنان لمخيلتنا وننظر إلى العواصف التي تتصاعد في الأفق.}: اختر منصة إعادة تدوير ذات سمعة جيدة، بعد الطلب يأتي مندوب SF للاستلام، وبعد الفحص يُحوّل المبلغ. احرص على اختيار\textbf{منصة تقدم "تقرير مسح البيانات"}واحتفظ بالإيصال.
  \item \textbf{التبرع العام}: يمكن التواصل مع مدارس الريف المحلية أو منظمات المجتمع المحلي الخيرية، الهواتف القديمة يمكن تعديلها لأجهزة لوحية تعليمية أو هواتف لكبار السن.\textbf{قبل التبرع نفّذ عملية المسح الكاملة، وأزل جميع بطاقات SIM/SD}.
  \item \textbf{الهواتف القابلة للطي}: تحويلها إلى جهاز ملاحة سيارة أو مركز تحكم للمنزل الذكي أو كاميرا مراقبة أو جهاز مراقبة طفل أو إطار صور إلكتروني. هاتف قديم يمكنه العمل لسنوات أخرى.
\end{itemize}

\begin{tcbinfo}{حقيقة النفايات الإلكترونية}
تحتوي لوحة دائرة الهاتف على\textbf{ذهب وفضة وبلاديوم ونحاس}وغيرها من المعادن الثمينة، وكل طن من اللوحات المستعملة يحوي ذهباً يعادل عشرات الأضعاف ما في خام الذهب. لكن إذا رُميت عشوائياً، فإن الرصاص والزئبق والكادميوم فيها تتسرب للتربة والمياه الجوفية، والتلوث يستمر لعقود.\textbf{الإعادة للتدوير الصحيحة تحمي البيئة وتتيح إعادة تدوير المعادن النادرة.}

\end{tcbinfo}

\begin{tcbpractice}{نظارات الواقع المعزز (AR)}
\begin{enumerate}
  \item آخر نسخة احتياطية كاملة: رفع الصور للسحابة ونسخ الملفات المهمة للحاسوب وتصدير سجلات الدردشة.
  \item تسجيل الخروج من جميع الحسابات: iCloud / Huawei / Xiaomi / Google، وإغلاق "Find My Phone".
  \item إخراج بطاقة SIM وSD card وفصل جميع أجهزة OTG.
  \item تنفيذ استعادة للمصنع مشفّرة، وتحديد "تهيئة التخزين" يدوياً عند الحاجة.
  \item (الأجهزة القديمة فقط) ملء التخزين بملفات غير مهمة ثم استعادة للمصنع مرة أخرى.
  \item اختيار المحطة التالية: الإعادة الرسمية بجهاز قديم أو منصة إعادة تدوير معتمدة أو تبرع عام أو إهداء للأسرة.
  \item تحديث "قائمة الميراث الرقمي"، وشطب الأجهزة المُستغنى عنها.
\end{enumerate}

\end{tcbpractice}

\begin{tcbpractice}{قائمة مهام هذا الشهر}
\begin{enumerate}
  \item تعيين "جهة اتصال الميراث" من Apple أو ضبط "مدير الحساب غير النشط" من Google لهاتفك الرئيسي.
  \item صياغة "ملحق معالجة الأصول الرقمية الشخصية" بالنموذج المرفق، والتواصل مع الأسرة حول مكان الحفظ.
  \item تحقق من وجود جهازين قديمين أو أكثر غير مستخدمين في المنزل، واحجز إعادة تدوير عبر منصة رسمية.
\end{enumerate}

\end{tcbpractice}
